\chapter{Experimental Setup}
\label{chap:experimental-setup}

This chapter defines what needs to be built and captured to execute the methodology of Chapter \ref{chap:methodology}. All items in this chapter are pending; none have been executed yet.

\section{Blender scene construction}

\todo[inline]{Pending: build the Blender scene. Required elements: (1) tether proxy geometry matching the real proxy's diameter and material appearance, (2) dark-enclosure interior geometry and matte-black material properties, (3) camera object with intrinsic parameters matching the calibrated real camera, (4) light source configured per Chapter \ref{chap:background}: Blackbody shader at approximately 5772K, Sun strength set to the real light source's measured irradiance for the laboratory-condition render, and to 1361 W/m\textsuperscript{2} (AM0) for the orbital-condition render.}

\subsection{Blender's own configuration reference for the Sun light}

Blender's own documentation, rather than a third-party tutorial, confirms two of the settings already used in the project's scene file. The Sun light object exposes an \texttt{energy} property described directly in the Blender Python API reference as ``Sunlight strength in watts per meter squared (W/m\textsuperscript{2})'', unbounded range, default 10.0 \cite{blender_sunlight_api}, which is the same irradiance unit discussed in Chapter \ref{chap:background} and confirms the value can be set directly with no conversion. The same object exposes an \texttt{angle} property, described as the ``Angular diameter of the Sun as seen from the Earth'', range 0 to $\pi$ radians, default 0.00918043 rad ($\approx$0.526\textdegree) \cite{blender_sunlight_api}. This matches, to three significant figures, the real solar angular diameter ($\approx$0.53\textdegree) already set on the Sun object in \texttt{SIMULATION/Simulation\_DarkBox.blend}, which is therefore confirmed correct against Blender's own default rather than only against the external astronomical figure used when it was first set.

Blender's manual is explicit that scene-linear light values in Cycles are unbounded (0 to infinity) and only compressed to a displayable image by the View Transform chosen in Color Management, not by the light values themselves \cite{blender_color_management}. The default View Transform in Blender 5.x, AgX, is described in the manual as offering approximately 16.5 stops of dynamic range and deliberately desaturating strongly overexposed colors rather than clipping them, so that bright regions roll off toward white the way a real camera's sensor does instead of producing flat, saturated color blocks \cite{blender_color_management}. This is the mechanism, native to Blender and not a custom workaround, that directly satisfies the requirement that the rendered light must not appear saturated: keeping View Transform set to AgX (the default) rather than Standard is the concrete action to take when the scene is rendered, and the Exposure slider (output equals the rendered value multiplied by $2^{\text{exposure}}$) is the correct control to fine-tune brightness afterward without reintroducing clipping.

\section{Laboratory photograph capture}

A reference stereo pair has been selected from the existing capture archive: the tether enters the frame from above the camera and falls toward the background, chosen because this configuration is simple to reproduce and reason about (no self-crossing, no ambiguous topology) for a first matched-pair test.

\begin{itemize}
    \item Left image: \texttt{data/captures/stereo\_20260213\_140535/left.jpg}
    \item Right image: \texttt{data/captures/stereo\_20260213\_140535/right.jpg}
    \item Metadata: \texttt{data/captures/stereo\_20260213\_140535/metadata.json} (capture resolution 1920$\times$1440, baseline 100 mm, synchronization quality "EXCELENTE", 10.4 ms time difference between cameras)
\end{itemize}

\todo[inline]{Pending: this pair was captured before the current calibration was invalidated (Chapter 7 recalibration item, Camera System Report), so it is usable as a 2D reference for mask/endpoint/path comparison but not for 3D reconstruction comparison. Also pending: document the physical pose of the tether proxy for this specific capture (approach angle, distance from camera, support/attachment point) so it can be replicated in the Blender scene.}

\section{Matched-pair generation}

\todo[inline]{Pending: for each real photograph captured above, render the corresponding Blender scene from the same camera position, orientation, and tether pose, producing one matched real/simulated pair per configuration tested. A minimum of 3 to 5 matched pairs, covering different tether poses (straight, curved, and one closer to a self-crossing configuration if feasible), is recommended before drawing conclusions from a single case.}

\section{Orbital-condition render generation}

\todo[inline]{Pending, and dependent on satisfactory agreement in the laboratory-condition comparison (Chapter \ref{chap:results}): reconfigure the validated Blender scene to represent orbital conditions (vacuum, no enclosure, AM0 solar illumination, unobstructed star-field or black background) and render the tether under representative deployment geometries.}
